% ============================================================
% 2. MARCO TEÓRICO
% ============================================================
\section{Marco Teórico}
\label{sec:marco-teorico}

\subsection{Máquinas virtuales y virtualización}
\label{sec:mt-virtualizacion}

La virtualización permite ejecutar múltiples sistemas operativos aislados
sobre un mismo hardware físico mediante un hipervisor, que intermedia el
acceso a los recursos de la máquina. En este trabajo se utiliza Oracle VM
VirtualBox, un hipervisor de tipo 2 (alojado sobre un sistema operativo
anfitrión) multiplataforma, disponible para GNU/Linux y Windows
\parencite{virtualbox}.

VirtualBox se administra tanto desde su interfaz gráfica como desde la
utilidad de línea de comandos \texttt{VBoxManage}, incluida en toda
instalación y disponible de forma idéntica en todos los sistemas
operativos soportados \parencite{virtualbox}. Esta característica es
relevante para la consola: cualquier operación que la aplicación deba
realizar sobre la máquina virtual (consultar su estado, obtener su
dirección IP, ejecutar un programa dentro de ella) puede expresarse como
llamadas a \texttt{VBoxManage}, lo que la convierte en la base de la
comunicación inicial con la máquina virtual.

Las Guest Additions son controladores y servicios que se instalan dentro
de la máquina virtual y mejoran la integración con el anfitrión:
redimensionado de pantalla, portapapeles compartido, directorios
compartidos y, en particular, la publicación de propiedades del sistema
invitado \parencite{guestadditions}. Entre estas propiedades se encuentra
la dirección IP de la máquina, publicada como una propiedad consultable
desde el anfitrión con el nombre
\texttt{/VirtualBox/GuestInfo/Net/0/V4/IP}. Este mecanismo resuelve un
problema práctico del proyecto: antes de que exista cualquier servicio de
red configurado dentro de la máquina (como SSH), la consola necesita
conocer su dirección para poder establecer la primera conexión.

\subsection{Sistema operativo GNU/Linux Debian}
% TODO: distribución Debian 13, estructura básica, gestión de archivos
% de configuración relevantes (/etc/issue, /etc/hosts).

\subsection{Gestión de paquetes con APT}
% TODO: funcionamiento de apt install / apt remove, repositorios.

\subsection{Protocolo y servicio SSH}
\label{sec:mt-ssh}

El protocolo Secure Shell (SSH) proporciona canales seguros sobre redes no
confiables, con autenticación cifrada, confidencialidad e integridad de
las comunicaciones \parencite{rfc4252}. En servidores GNU/Linux el
servidor de referencia es OpenSSH, que implementa el servicio
\texttt{sshd} y el cliente homónimo.

SSH admite autenticación por contraseña y por par de llaves
criptográficas asimétricas. En el esquema por llaves, el cliente posee
una llave privada y el servidor guarda la llave pública correspondiente en
el archivo \texttt{\textasciitilde/.ssh/authorized\_keys} de cada usuario;
el servidor verifica la firma del cliente sin transmitir secretos, y es
posible deshabilitar por completo el acceso por contraseña
\parencite{rfc4252}. Este es el mecanismo que la consola configura en la
máquina virtual: una vez instalado y habilitado SSH, toda la comunicación
posterior entre la consola y la máquina se realiza por este canal,
reemplazando la comunicación inicial basada en Guest Additions.

Además de la ejecución interactiva de comandos, el cliente OpenSSH
permite ejecutar una orden remota de forma no interactiva (\texttt{ssh
usuario@host orden}) y transferir archivos mediante \texttt{scp}, lo que
lo convierte en una interfaz de programación completa para automatizar
administración de sistemas \parencite{donovan}.

\subsection{Gestión de servicios con systemd}
% TODO: unidades de systemd, ciclo de vida (start, stop, restart,
% status), habilitar/deshabilitar al arranque, archivos .service.

\subsection{Balanceo de carga con HAProxy}
% TODO: concepto de balanceo de carga, front-end y back-end,
% algoritmos de distribución, configuración básica de HAProxy.

\subsection{Lenguaje de programación Go y aplicaciones web}
\label{sec:mt-go}

Go es un lenguaje compilado, de tipado estático y concurrencia ligera,
diseñado para producir binarios simples y multiplataforma
\parencite{donovan}. Compilar el mismo código para GNU/Linux y Windows
sin dependencias en tiempo de ejecución encaja con el objetivo de esta
consola, ya que la aplicación debe poder ejecutarse en el equipo
anfitrión independientemente de su sistema operativo.

La interfaz web se sirve con el paquete estándar \texttt{net/http} y
sigue una arquitectura cliente-servidor: el navegador presenta los
módulos de administración y envía formularios; el servidor en Go valida
las entradas, ejecuta las operaciones sobre la máquina virtual y devuelve
el resultado.

Para ejecutar operaciones externas, Go ofrece el paquete \texttt{os/exec},
que crea y controla procesos del sistema operativo capturando su salida,
su error y su código de salida \parencite{donovan}. Este es el primitivo
sobre el que se construye la biblioteca de conectores: la consola no
invoca \texttt{vboxmanage} ni \texttt{ssh} directamente desde su lógica,
sino a través de procesos externos lanzados de forma controlada, lo que
permite usar los mismos binarios en Windows y en GNU/Linux.

\subsection{Arquitectura de puertos y adaptadores}
\label{sec:mt-puertos-adaptadores}

Para mantener el núcleo de la aplicación independiente de la
infraestructura concreta se aplica el patrón de arquitectura de puertos
y adaptadores (también llamado arquitectura hexagonal)
\parencite{cockburn}. En este patrón, el núcleo define \emph{puertos}:
interfaces que expresan lo que la aplicación necesita del exterior
(ejecutar un programa, abrir una sesión con la máquina, conocer su
dirección). Los \emph{adaptadores} son implementaciones concretas de esos
puertos, que traducen el lenguaje de una tecnología específica al
lenguaje del dominio.

En este proyecto los puertos son \texttt{Executor} (ejecutar un proceso
en el anfitrión), \texttt{Machine} (obtener dirección y estado de la
máquina virtual), \texttt{Connector} (establecer sesiones) y
\texttt{Session} (ejecutar órdenes y transferir archivos dentro de la
máquina). Los adaptadores implementados son \texttt{localproc} (procesos
del sistema operativo), \texttt{vbox} (VirtualBox mediante
\texttt{vboxmanage} y Guest Additions) y \texttt{ssh} (cliente OpenSSH).
De este modo, el núcleo no sabe que se usa Guest Additions ni VirtualBox:
solo conoce interfaces, y cambiar una tecnología por otra (u otra
distinta) no afecta a la lógica de la consola.
